\subsection{The actual tail algebra for one independent degree-four class}

\paragraph{A-S09-21 (Proved; a nonfree coefficient-tail attachment and its exact map data).}
Let $X$ be a free identity-disk algebra as in A-S09-18,
and choose one odd form-one generator $x$ of weight $s$,
with $dx=y$.  Put
\[
 z_0=\gamma_2(y),\qquad b_0=xy,\qquad db_0=2z_0,
 \qquad L=I_N\langle2s\rangle,\quad N\geq1.          \tag{S09.92}
\]
Let $F$ be the free relative graded PD algebra over
$X$ on two copies of the actual graded $S$-module $L$,
in form degrees two and three.  Denote the generators
by $v(l),u(l)$, with their $S$-linear relations, and set
\[
 du(l)=l z_0,\qquad dv(l)=2u(l)-l b_0,
 \qquad Y=\operatorname{im}(F\to F\otimes\mathbb Q).
 \tag{S09.93}
\]
The notation $l z_0$ uses the underlying unshifted
coefficient $l\in I_N$; the shift supplies its weight.
These formulas define an actual strict PD differential
algebra $Y$ containing $X$, with torsion-free terms,
unchanged weights less than $N+2s$, and
\[
 H^0(Y)=S,\qquad H^1(Y)=H^2(Y)=H^3(Y)=0,
 \qquad
 H^4(Y)=H^4(X)\big/\bigl((I_N/2I_N)[z_0]\bigr).
 \tag{S09.94}
\]
The statement is integral in every weight.

For every torsion-free strict PD target $A$ with a
specified map $f:X\to A$, the maps $Y\to A$ extending
$f$ are in exact bijection with pairs of graded
$S$-linear maps
\[
 U:L\to A^3,\qquad V:L\to A^2,
 \qquad dU(l)=l f(z_0),\qquad
 dV(l)=2U(l)-l f(b_0).                               \tag{S09.95}
\]
This identifies the actual target-map data required
by this tail algebra.  It does not supply solutions
of (S09.95) from eventual exactness alone.

\begin{proof}
Both right sides of (S09.93) are $S$-linear in $l$,
so they respect the defining linear relations.
The equality $db_0=2z_0$ gives $d^2v=0$; the
closedness of $z_0$ gives $d^2u=0$.
The special derivation rule extends these assignments
over the free PD presentation.  Its torsion ideal
is preserved by differential and divided powers,
so the displayed torsion-free image is a strict
PD quotient.  The graded augmentation of the
free factor splits $X\to F$.  Since $X$ is
torsion-free, this inclusion remains injective
after passing to $Y$.  The weights and forms of
the generators prove the assertions about lower
weights and degrees zero and one.

Here the linear new-generator terms in degrees
two through four are exactly the relevant terms
of $X\otimes_S L$; they have no integer torsion.
Indeed $X$ is free over $S$ on its PD and exterior
monomial basis, while $L$ is free as an abelian
group.  The explicit identity-disk decomposition
also proves
\[
 H^0(X\otimes_S L)=L,
 \qquad H^j(X\otimes_S L)=0\quad(1\leq j\leq3).
 \tag{S09.96}
\]
Every unit pair contracts after tensoring, and
the first remaining pair is multiplication by
two in degrees three and four.  It remains
injective on $L$, which is torsion-free.
Thus the linear new-generator terms are unchanged
by the torsion quotient in (S09.93).

The degree-two and degree-three proof of A-S09-20
now applies with (S09.96) in place of its free
coefficient calculation.  In particular the
degree-three group is the kernel of the actual
map
\[
                      L/2L\longrightarrow H^4(X),
                                     \quad l\longmapsto l[z_0].
 \tag{S09.97}
\]
By (S09.80), $[z_0]$ generates a direct summand
$S/2S$ of $H^4(X)$.  The coefficient inclusion
$I_N/2I_N\to S/2S$ is injective: its quotient
$S/I_N$ is free as an abelian group, or this can
be checked separately in each weight.  Hence
(S09.97) is injective and $H^3(Y)=0$.

For degree four the only additional new terms
are the quadratic divided powers of $v(L)$.
After rationalization, $L_{\mathbb Q}$ is a free
rank-one graded $S_{\mathbb Q}$-module, generated
by $T^{[N]}$ with the indicated shift.  The part
of the differential from quadratic $v$ terms
to the distinct $uv$ terms is therefore the
injective derivative with coefficient two in
the ordinary rational polynomial algebra.
A closed quadratic contribution has zero
rationalization; it is zero in the torsion-free
algebra $Y$.  For the remaining linear terms,
(S09.96) removes the $v$ coefficients and then
the $u$ coefficients exactly as in A-S09-20.
Consequently every closed four-form is represented
by an old one, and the old boundaries newly
created are precisely the span in (S09.97).
This proves (S09.94).

The rational inclusion is also an actual
quasi-isomorphism.  On the free rank-one
$L_{\mathbb Q}$ generator $l_0=T^{[N]}$, replace
$u(l_0)$ by $u(l_0)-l_0b_0/2$.
Its differential becomes zero and that of
$v(l_0)$ becomes twice this new generator.
This is an invertible PD change of variables
over $\mathbb Q$, and the resulting even-to-odd
disk contracts over $X_{\mathbb Q}$.

Finally the free PD universal property identifies
maps from $F$ over $X$ with the generator values
in (S09.95); the two displayed equations are
exactly the requirement to commute with the
differential.  Every map from $F$ to torsion-free
$A$ kills the torsion ideal, so it factors uniquely
through $Y$.  Conversely a map from $Y$ gives
exactly those generator values.  This proves
the asserted bijection of actual morphisms.
\end{proof}

This construction retains the nonfree coefficient
module $I_N$ and its actual linear relations.  It
does not replace it by a monomial-tail enlargement,
nor infer the missing maps (S09.95) from a derived
Hom computation.  Those particular maps remain
the necessary further input for this tail step
in the common-source construction.
